\documentclass[12pt, a4paper]{article}

\usepackage[T1]{fontenc}
\usepackage[utf8]{inputenc}
\usepackage[spanish, es-tabla]{babel}
\usepackage[a4paper, left=2.4cm, right=2.4cm, top=2.2cm, bottom=2.2cm]{geometry}
\usepackage{xcolor}
\usepackage{graphicx}
\usepackage{booktabs}
\usepackage{tabularx}
\usepackage{enumitem}
\usepackage{fancyhdr}
\usepackage[hidelinks]{hyperref}

\definecolor{azul}{HTML}{2F5BD8}
\definecolor{grafito}{HTML}{22252A}
\definecolor{gris}{HTML}{6B7280}
\definecolor{grisclaro}{HTML}{F0F1F4}

\renewcommand{\familydefault}{\sfdefault}
\color{grafito}
\linespread{1.12}\selectfont
\setlength{\parindent}{0pt}
\setlength{\parskip}{0.55em}

\pagestyle{fancy}
\fancyhf{}
\fancyhead[L]{\footnotesize\color{gris} Guion de la defensa · Brain Viz}
\fancyhead[R]{\footnotesize\color{gris} Pablo Moliz Arias}
\fancyfoot[C]{\footnotesize\color{gris}\thepage}
\renewcommand{\headrulewidth}{0.4pt}

% Cabecera de cada diapositiva: número, título, duración y tiempo acumulado.
\newcommand{\dia}[4]{%
	\vspace{1.1em}%
	\noindent\colorbox{grisclaro}{%
		\begin{minipage}{\dimexpr\textwidth-2\fboxsep}
			\vspace{0.25em}
			{\color{azul}\bfseries Página #1}\quad{\bfseries #2}\hfill
			{\footnotesize\color{gris} #3 · hasta #4}
			\vspace{0.25em}
		\end{minipage}}%
	\vspace{0.5em}\par}

% Indicaciones de lo que hay que hacer, no de lo que hay que decir.
\newcommand{\accion}[1]{%
	\vspace{0.2em}\par\noindent
	{\color{azul}\rule{2.5pt}{1em}}\hspace{0.5em}{\itshape\color{azul} #1}\par
	\vspace{0.2em}}

\newcommand{\pausa}{{\color{azul}\bfseries\,/\,}}

\begin{document}

% ---------------------------------------------------------------- portada ---
\thispagestyle{empty}
\vspace*{2.5cm}
\begin{center}
	{\footnotesize\color{gris} TRABAJO FIN DE GRADO}\\[1.2em]
	{\LARGE\bfseries Guion de la defensa}\\[0.6em]
	{\color{azul}\rule{4em}{2.5pt}}\\[0.8em]
	{\large Aplicación web para la visualización 3D de redes cerebrales}\\[0.3em]
	{\large\bfseries Brain Viz}\\[2.2em]
	Pablo Moliz Arias\\[0.3em]
	{\color{gris} Tutor: Juan Ruiz de Miras}\\[2.2em]
	{\footnotesize\color{gris} Duración prevista: 20 minutos\\
	 50 páginas: 37 de contenido, 5 separadores, 1 de vídeo y 7 de reserva}
\end{center}
\vfill
\newpage

% -------------------------------------------------------- cómo usar esto ---
\section*{Antes de empezar}

\textbf{Cómo está montado este guion.} Cada apartado corresponde a una diapositiva. Arriba tienes su
número de página, cuánto debería durar y en qué minuto tendrías que ir al terminarla. El texto está escrito para
decirlo, no para leerlo palabra por palabra: apréndete la idea y las cifras, y suéltalo con tus
palabras. Las líneas en azul no se dicen, son cosas que hacer.

\textbf{El símbolo \pausa} marca dónde conviene respirar y dejar un silencio corto. Son los momentos
en los que se entiende lo que acabas de decir.

\textbf{Los números son los de página del PDF}, que es lo que ves en el visor y en el mando. Entre
bloque y bloque hay una página separadora con el nombre de la sección, que pasas sin decir nada.

\textbf{Ritmo.} Hay muchas diapositivas y cada una dura poco, entre veinte y cuarenta y cinco
segundos. La idea es que vayas cambiando a buen paso: cada una tiene una sola cosa que contar, y
cuando la hayas contado, pasas. No te quedes parado en ninguna.

\textbf{No leas la pantalla.} Las diapositivas llevan a propósito muy poco texto, casi siempre una
frase o cuatro palabras. Eso que se ve es el titular; lo que explica ese titular lo dices tú. Si te
descubres leyendo en voz alta lo que hay escrito, párate y cuéntalo con tus palabras.

\textbf{Si vas con retraso}, recorta las páginas 4, 15 y 27, que son las que menos se echan de
menos. Lo que no se recorta nunca es la 22 (la idea que resolvió el rendimiento) ni la 28 y la 29
(cómo validé las métricas y el error que encontré): son las que enseñan que hay trabajo de
ingeniería detrás.

\vspace{0.5em}
\textbf{Lista de comprobación del día:}
\begin{itemize}[leftmargin=1.4em, itemsep=0.2em]
	\item \textbf{Abre el vídeo antes de empezar}, en un reproductor aparte, en pausa y a pantalla
		completa. Así al llegar a la página 38 solo tienes que cambiar de ventana y darle.
	\item Comprueba que el vídeo \textbf{no tiene volumen}. No lleva audio, pero por si acaso.
	\item Lleva el PDF y el vídeo en el portátil, en un USB y en tu correo.
	\item \textbf{Nada de demostración en directo.} Si el tribunal pide ver la aplicación de verdad,
		la tienes en \texttt{brain-viz.onrender.com}, pero eso es para el turno de preguntas.
	\item Ten a mano el PDF de la memoria por si preguntan por una tabla concreta.
\end{itemize}

\vspace{0.6em}
\textbf{Reparto del tiempo}
\vspace{0.3em}

\begin{center}\footnotesize
\begin{tabularx}{\textwidth}{@{}l X r r@{}}
	\toprule
	\textbf{Bloque} & \textbf{Páginas} & \textbf{Dura} & \textbf{Acaba en} \\
	\midrule
	Apertura              & 1        & 0:20 & 0:20 \\
	El problema           & 3 a 7    & 2:35 & 2:55 \\
	Qué se pedía          & 9 a 12   & 1:45 & 4:40 \\
	Cómo lo he construido & 14 a 34  & 10:40 & 15:20 \\
	Cómo sé que funciona  & 36 y 37  & 0:45 & 16:05 \\
	\textbf{Vídeo}        & 38       & \textbf{2:30} & 18:35 \\
	Conclusiones          & 40 a 43  & 1:05 & 19:40 \\
	\bottomrule
\end{tabularx}
\end{center}

\newpage

% ================================================================== guion ===
\section*{El guion}

\dia{1}{Portada}{0:20}{0:20}

Buenos días. Soy Pablo Moliz Arias y vengo a presentar mi Trabajo Fin de Grado, dirigido por Juan
Ruiz de Miras. Se titula «Aplicación web para la visualización 3D de redes cerebrales», y la
aplicación se llama Brain Viz. \pausa

\accion{Cambia de diapositiva mientras dices la última frase. La página 2 es un separador: pasas sin
decir nada.}

\dia{3}{El cerebro es una red}{0:35}{0:55}

El punto de partida es una idea sencilla. El cerebro se puede representar como una red: cada región
es un nodo, colocado en las coordenadas que le corresponden anatómicamente, y cada conexión entre
dos regiones es una arista con un peso, que mide cómo de intensa es esa conexión. El ejemplo que uso
en todo el trabajo tiene noventa regiones. \pausa

\dia{4}{Y sobre una red se puede medir}{0:30}{1:25}

Lo interesante de mirarlo así es que la teoría de grafos ya tiene medidas hechas para responder a
preguntas que en neurociencia importan. Qué regiones actúan como puntos de paso del resto. Si la red
está agrupada en bloques muy conectados entre sí. O si se separa en partes que no se comunican. Son
preguntas que sobre una imagen del cerebro no se pueden contestar, y sobre un grafo sí.

\dia{5}{Lo que ya existe: en el escritorio}{0:35}{2:00}

Antes de programar nada revisé qué había, y hay herramientas buenas, pero ninguna encajaba del todo.

BrainNet Viewer es probablemente la más usada, y de hecho de ella he tomado el formato de ficheros,
pero depende de MATLAB, que es comercial. Connectome Workbench es la referencia del Human Connectome
Project, muy completa, pero con instalación y formatos propios. Y Gephi es potentísima para grafos,
aunque no coloca los nodos donde están anatómicamente, que aquí es justo lo que importa.

\dia{6}{Lo que ya existe: en el navegador}{0:35}{2:35}

En el navegador, Vizaj es la más cercana a lo que yo quería, aunque se ejecuta entera en el cliente
y no tiene ningún análisis en servidor. BioImage Suite Web sí es una suite completa, pero el visor
de conectividad es uno más de sus muchos módulos, con la complejidad que eso trae. Y Neuroglancer,
que es el visor web más conocido, está pensado para volúmenes de imagen, no para redes. \pausa

\dia{7}{El hueco}{0:20}{2:55}

Así que el hueco estaba bastante claro.

\accion{Deja que lean la frase un segundo antes de seguir. No la leas tú en voz alta.}

\dia{9}{Objetivo}{0:20}{3:15}

De ahí sale el objetivo general del trabajo.

\accion{Igual: pausa, que lo lean, y sigues.}

\dia{10}{Cinco objetivos}{0:35}{3:50}

Y cinco objetivos específicos. Los tres primeros son de funcionalidad: ver la red en 3D de forma
interactiva, poder analizarla con métricas de grafos, y poder navegar, filtrar y seleccionar sobre
ella. Los dos últimos son de ingeniería: que el sistema sea mantenible, con el procesamiento
separado de la visualización, y que se pueda desplegar. Volveré a estos cinco al final.

\dia{11}{Requisitos}{0:20}{4:10}

Eso se concretó en dieciocho requisitos funcionales, siete no funcionales y cuatro de información, y
en tres roles: visitante, investigador y administrador.

\dia{12}{Cuatro sprints}{0:30}{4:40}

Los requisitos se priorizaron en cuatro niveles, y cada nivel es un sprint. El primero deja la red
en pantalla, el segundo permite explorarla, el tercero la analiza y el cuarto añade los usuarios.
Cada uno acaba en algo que se puede enseñar y probar, no en media funcionalidad. Y entre sprint y
sprint hubo revisión con mi tutor, que como verán cambió varias cosas. \pausa

\dia{14}{Arquitectura}{0:40}{5:20}

La arquitectura es cliente-servidor desacoplada. El navegador dibuja y el servidor procesa, y hablan
por una API.

La regla que me impuse desde el principio, y que se mantuvo los cuatro sprints, es que el servidor
no sabe absolutamente nada de cómo se dibuja la red, y el cliente no calcula ninguna métrica. Eso
permite cambiar cualquiera de los dos lados sin tocar el otro, y es lo que hizo que el cuarto
sprint, que fue el que más cambió, no rompiera nada de lo anterior.

\dia{15}{Tecnologías}{0:35}{5:55}

Sobre las tecnologías, más que la lista me interesa el porqué. Flask porque la API es pequeña y de
lectura, y un marco completo como Django habría traído mucho que no iba a usar. NetworkX porque
implementa las métricas del área y está contrastada. React porque la interfaz es estado que cambia
continuamente. Y Three.js a través de react-three-fiber, que me permite describir la escena 3D igual
que el resto de la interfaz.

\dia{16}{Sprint 1: la red en pantalla}{0:35}{6:30}

El primer sprint deja la red en pantalla. Se leen los ficheros del formato de BrainNet Viewer, y
elegí ese precisamente porque es el que ya tiene la gente del área: si hubiera inventado un formato
propio, cualquiera que quisiera probarla tendría que convertir sus datos antes. Los noventa nodos se
dibujan en sus coordenadas anatómicas reales y la cámara responde sin esperas.

\dia{17}{Umbral 0,00}{0:20}{6:50}

El segundo sprint convierte esa imagen en algo explorable. Esto es la red entera, con sus
cuatrocientas once conexiones, que es prácticamente una maraña.

\dia{18}{Umbral 0,70}{0:25}{7:15}

Y esto es lo mismo subiendo el umbral, dejando solo las conexiones más intensas. Ahí aparece la
estructura: se ve qué zonas están fuertemente conectadas entre sí. \pausa

\dia{19}{Seleccionar una región}{0:25}{7:40}

Al seleccionar una región se resaltan sus conexiones y el panel muestra sus datos. Esta es el
putamen izquierdo, con grado once de once: sus once conexiones están todas visibles.

\dia{20}{El grado cambia con el umbral}{0:30}{8:10}

Y aquí está una de las correcciones de mi tutor. Si ahora subo el umbral, el mismo nodo pasa a
mostrar siete de once. Antes yo enseñaba siempre el grado en la red completa, y al filtrar eso deja
de ser lo que la persona está viendo. Parece un detalle y cambia por completo la lectura.

\dia{21}{El problema de rendimiento}{0:30}{8:40}

Ahora la parte que más me interesa contarles, porque es donde hubo un problema de verdad.

Mi primera versión del filtrado recorría todas las aristas y reconstruía la geometría entera cada
vez que se movía el control. Con la red pequeña colaba, pero con miles de conexiones se notaba en la
mano: el control iba a tirones. \pausa

\dia{22}{La propiedad que lo simplifica}{0:40}{9:20}

La solución no fue optimizar ese código, sino darme cuenta de una propiedad de los datos. Si el
servidor envía las aristas ya ordenadas por peso, entonces las visibles para cualquier umbral son
siempre las primeras de la lista. Nunca son un conjunto disperso.

Y eso lo cambia todo. La geometría se construye una sola vez, al cargar, y filtrar pasa a ser
simplemente decirle a la tarjeta gráfica cuántas de esas aristas tiene que dibujar.

\dia{23}{El resultado}{0:25}{9:45}

Cuántas son se averigua con una búsqueda binaria sobre la lista ordenada. El resultado es que mover
el control deja de depender del tamaño de la red. \pausa

\accion{Si preguntan: comprobé la búsqueda binaria contra un recuento lineal en 207 casos, sin
ninguna discrepancia.}

\dia{24}{Lo que salió de las revisiones}{0:40}{10:25}

Quiero detenerme un momento en las revisiones, porque cambiaron el resultado más de lo que yo
esperaba.

Mi tutor me pidió que las conexiones se vieran más gruesas. Parecía cuestión de cambiar un número, y
resultó que la especificación de WebGL no obliga a los navegadores a dibujar líneas de más de un
píxel, y en la práctica ignoran ese parámetro. Hubo que cambiar por completo la forma de dibujarlas.

El segundo ya lo han visto, el grado con el umbral. Y el tercero, que las conexiones del nodo
elegido se distingan manteniendo su color, en lugar de teñirse.

Ninguna de las tres se me habría ocurrido probando yo solo la aplicación, precisamente porque yo ya
sabía lo que hacía cada cosa. \pausa

\dia{25}{Las métricas de la red}{0:25}{10:50}

El tercer sprint es el análisis. Se calculan cuatro métricas globales, que describen la red entera,
y tres medidas de centralidad, que dan un valor por región: grado, intermediación y cercanía.

\dia{26}{El mapeo visual}{0:30}{11:20}

Pero un número en una tabla, sobre noventa regiones, no dice gran cosa. Por eso lo más útil de este
sprint es el mapeo visual: el tamaño y el color de cada nodo pasan a representar la métrica elegida.
Esto que ven es la intermediación, y de un vistazo se identifica qué regiones son los puntos de paso
de la red. La respuesta está en la escena, no en una tabla.

\dia{27}{El buscador}{0:25}{11:45}

Se añadió también un buscador. Y tiene una historia: yo lo daba por terminado porque el filtrado
funcionaba, hasta que me di cuenta de que las etiquetas del atlas son abreviaturas, y quien escribe
«hipocampo» no obtiene nada. Lo que lo arregló no fue tocar el filtrado, sino desplegar la lista de
las noventa regiones para poder elegir directamente.

\dia{28}{¿Y esos números están bien?}{0:35}{12:20}

Y aquí viene lo que para mí es la parte más importante del trabajo. Si alguien va a usar estos
números para investigar, tienen que estar bien.

Así que no me limité a comprobar que salían números: contrasté las cinco métricas con
implementaciones escritas por mí sin usar la librería, incluida una del algoritmo de Brandes para la
intermediación, que no es precisamente trivial. \pausa

\dia{29}{Y apareció un error real}{0:35}{12:55}

Y esa comprobación encontró un error de verdad. Verificando que la suma de los grados fuera el doble
del número de aristas, que es una propiedad que se cumple siempre, descubrí que todos los nodos
tenían exactamente el mismo grado. La librería estaba creando una conexión por cada celda de la
matriz, incluidas las que valen cero.

Lo relevante es que ese fallo no daba ningún síntoma en pantalla. La aplicación se veía perfecta.
Sin esa comprobación habría llegado hasta el final del proyecto. \pausa

\dia{30}{Sprint 4: cuentas y catálogo}{0:30}{13:25}

El cuarto sprint convierte la herramienta en un sistema con usuarios. Aparecen las cuentas, el
catálogo de redes, la exportación de resultados y un panel de administración.

\dia{31}{Redes propias}{0:20}{13:45}

Y cada investigador puede subir sus propias redes, que solo ve él.

\dia{32}{Lo que hubo que resolver}{0:35}{14:20}

Este fue con diferencia el sprint que más impacto tuvo en la arquitectura, porque los tres primeros
compartían algo que simplificaba mucho las cosas: el sistema no guardaba nada.

Con los datos privados aparecen tres preguntas nuevas. Las contraseñas no se guardan, se guarda un
resumen del que no se puede volver atrás. La sesión va firmada, en una cookie que el navegador
impide leer desde el código. Y cada consulta filtra por el usuario de la sesión, sin fiarse nunca de
lo que pida el cliente. Fue lo primero que cubrí con pruebas automáticas. \pausa

\dia{33}{En el móvil}{0:25}{14:45}

Uno de los requisitos no funcionales es que la interfaz se adapte al tamaño de la pantalla. En un
teléfono no caben dos columnas de paneles, así que se cambia cuándo se muestran, no cuánto ocupan:
se abren desde una barra inferior. Y el criterio no puede ser solo el ancho, porque un teléfono
girado, como este, tiene ancho de sobra y casi nada de altura.

\dia{34}{Despliegue}{0:35}{15:20}

Sobre el despliegue tengo que contar algo que no tenía previsto. Durante los tres primeros sprints
mi despliegue era el entorno de desarrollo con contenedores, y yo lo daba por bueno. Al ir a
publicarlo descubrí que no valía: eran dos servicios que se hablaban por un nombre interno de la red
de Docker, que fuera de ella no existe.

Tuve que preparar una imagen distinta, en dos etapas, con todo el sistema en un único contenedor. La
lección es que «funciona en mi máquina» y «se puede desplegar» no son lo mismo. \pausa

\dia{36}{Validación}{0:25}{15:45}

Cómo he comprobado que esto funciona: diecisiete pruebas automáticas que se lanzan con una sola
orden, treinta y dos pruebas de aceptación, una por cada requisito, y las cinco métricas
contrastadas con código propio.

\dia{37}{Rendimiento}{0:20}{16:05}

Y sobre rendimiento, multiplicar por treinta el número de conexiones no produjo degradación
apreciable. Lo que sí dejo señalado en la memoria, porque prefiero decirlo a darlo por bueno, es que
la medida de los sesenta fotogramas por segundo no la considero concluyente: el mismo resultado
salía midiendo una página vacía.

\dia{38}{Vídeo de la aplicación}{2:30}{18:35}

\accion{Cambia a la ventana del reproductor, que ya tiene que estar abierta y en pausa, y dale al
play. El vídeo no tiene sonido: lo vas comentando tú. Estos son los momentos, con su minuto.}

Y ahora, en lugar de una demostración en directo, les enseño un vídeo de la aplicación funcionando.

\accion{0:00 · La red de ejemplo, noventa regiones en sus coordenadas anatómicas.}

Esta es la red de ejemplo. Cada esfera es una región y el color indica a qué grupo anatómico
pertenece. Arrastrando se gira la escena.

\accion{0:25 · Empieza a girar y luego acerca la cámara.}

La cámara responde sin esperas, y con la rueda se acerca y se aleja.

\accion{0:50 · Sube el control de umbral, despacio.}

Ahora estoy subiendo el umbral. Fíjense en el contador de la izquierda: va bajando de cuatrocientas
once conexiones, y a la vez la maraña se deshace y aparece la estructura de la red.

\accion{1:15 · Vuelve a bajar el umbral y selecciona un nodo.}

Vuelvo a bajarlo, y al tocar una región se resaltan sus conexiones y el panel de la derecha muestra
sus datos: su grupo, su grado y sus dos medidas de centralidad.

\accion{1:35 · Cambia el criterio de color a intermediación.}

Aquí cambio el criterio de color a intermediación. El tamaño y el color de cada nodo pasan a
representar esa métrica, y los puntos de paso de la red se ven inmediatamente.

\accion{1:55 · Buscador y catálogo de redes.}

Este es el buscador, con la lista de las noventa regiones. Y este el catálogo, donde están las redes
públicas.

\accion{2:10 · Inicio de sesión y redes propias.}

Y por último el acceso con cuenta. Al entrar como investigadora aparecen sus propias redes, que
nadie más ve.

\accion{Cuando acabe, vuelve al PDF. La página 39 es un separador.}

\dia{40}{Los cinco objetivos}{0:15}{18:50}

Volviendo a los cinco objetivos del principio: los cinco están cumplidos, y los dieciocho requisitos
funcionales están implementados.

\dia{41}{Lo que me llevo}{0:30}{19:20}

Si tuviera que quedarme con cuatro cosas: que buscar la propiedad de los datos rinde mucho más que
optimizar el código; que las pruebas no son un trámite, porque el error del grado no se veía en
pantalla; que «terminado» no es lo mismo que «funciona», como pasó con el buscador; y el valor de
que te revisen.

\dia{42}{Trabajo futuro}{0:20}{19:40}

Quedan líneas abiertas: detectar comunidades y comparar dos redes entre sí, admitir más formatos y
más atlas, poder compartir una red con personas concretas, y sobre todo la accesibilidad, que es la
carencia más clara porque la aplicación depende mucho del color.

\dia{43}{Gracias}{0:10}{19:50}

Y hasta aquí. Muchas gracias por su atención, quedo a su disposición para las preguntas.

\accion{Deja esta diapositiva en pantalla durante las preguntas.}

\newpage

% =============================================================== preguntas ===
\section*{Preguntas que pueden hacerte}

Están ordenadas de más a menos probable. Respuestas cortas: contesta y para, no sigas hablando por
llenar el silencio.

\vspace{0.5em}

\textbf{¿Por qué Flask y no Django o FastAPI?}\\
Porque la API es pequeña y básicamente de lectura. Django trae un modelo completo, panel de
administración y muchas cosas que no iba a usar. FastAPI habría encajado, pero Flask me daba lo que
necesitaba con menos piezas y con más documentación disponible. En la memoria está la comparación
completa.

\textbf{¿Cómo escala con redes grandes?}\\
El dibujado aguanta bien: probé con quinientos nodos y más de doce mil conexiones y no hubo
degradación apreciable, gracias a que la geometría no se reconstruye. Lo que sí se nota es el
cálculo de la intermediación, que pasa de doce milisegundos a algo más de medio segundo. Guardar el
resultado ya procesado lo resuelve a partir de la segunda consulta, pero la primera sigue costando
lo mismo. Para redes bastante mayores lo correcto sería calcularlo en segundo plano y avisar cuando
esté listo.

\textbf{¿Cómo garantizas que un usuario no ve las redes de otro?}\\
Cada red guarda a qué usuario pertenece, y toda consulta filtra por el usuario de la sesión antes de
devolver nada. No se confía en lo que pida el cliente. Y es lo primero que cubrí con pruebas
automáticas: hay pruebas que intentan leer y borrar la red de otro usuario pidiéndola directamente
por su identificador, y comprueban que el sistema responde que no existe.

\textbf{¿Dónde se guardan las contraseñas?}\\
No se guardan. Se guarda un resumen criptográfico con sal, del que no se puede recuperar la
contraseña original. Al iniciar sesión se compara el resumen, nunca el texto.

\textbf{¿Por qué sesión con cookie y no un token JWT?}\\
Porque cliente y servidor se sirven desde el mismo sitio, así que la cookie es más simple y más
segura por defecto: va firmada, el navegador impide leerla desde el código y se puede invalidar en
el servidor, que con un token es justo lo complicado. Un token tendría sentido si hubiera clientes
de terceros consumiendo la API.

\textbf{¿Cómo validaste las métricas?}\\
Escribí mi propia implementación de las cinco, sin usar la librería, incluida una del algoritmo de
Brandes para la intermediación, y comparé los resultados. Además comprobé propiedades que se cumplen
siempre, como que la suma de los grados sea el doble del número de aristas. Fue precisamente esa
comprobación la que destapó que la librería estaba creando aristas de peso cero.

\textbf{¿Qué es exactamente la intermediación y para qué sirve aquí?}\\
Mide cuántos de los caminos más cortos entre pares de regiones pasan por una región concreta. Una
región con intermediación alta es un punto de paso: si se daña, la comunicación entre partes de la
red se resiente. Por eso en neurociencia interesa identificarlas.

\textbf{¿Se cumplen los sesenta fotogramas por segundo del requisito?}\\
No lo puedo afirmar con la medida que hice. Medí con dos redes de tamaños muy distintos y salió lo
mismo en las dos, pero también salía lo mismo midiendo una página vacía, lo que indica que el techo
lo ponía el entorno de medida. Lo que sí puedo afirmar es que multiplicar por treinta las conexiones
no produjo degradación apreciable. Prefiero dejarlo señalado como pendiente antes que darlo por
bueno.

\textbf{¿Por qué el atlas AAL90? ¿Sirve para otros?}\\
Lo uso como ejemplo porque es muy común y los datos están disponibles, pero la aplicación no depende
de él: lee cualquier par de ficheros en el formato de BrainNet Viewer, con las coordenadas y la
matriz. De hecho cualquier usuario puede subir los suyos.

\textbf{¿Qué diferencia real hay con Vizaj, que es lo más parecido?}\\
Vizaj se ejecuta entero en el navegador y se centra en la representación. Brain Viz separa el
análisis en un servicio de servidor, que es lo que permite calcular métricas de grafos sin cargar al
navegador, guardar resultados y tener usuarios con datos propios. Y está pensada para desplegarse de
forma reproducible.

\textbf{¿Por qué no hay pruebas automáticas de la interfaz?}\\
Es una carencia que asumo. Las pruebas automáticas cubren el servidor, que es donde están los datos
y las reglas de acceso, que me pareció lo prioritario. La interfaz la validé con las treinta y dos
pruebas de aceptación, que son manuales y están documentadas una a una en la memoria. Con más tiempo
habría añadido pruebas de extremo a extremo.

\textbf{¿Qué pasa si el fichero que sube el usuario está mal?}\\
Se valida antes de aceptarlo: que los dos ficheros vayan juntos, que la matriz sea cuadrada y que
sus dimensiones coincidan con el número de nodos. Si algo no cuadra se rechaza con un mensaje que
dice qué falla, en lugar de guardarlo y fallar después al dibujarlo.

\textbf{¿Por qué Render y no un servidor de la universidad?}\\
Porque quería que el resultado fuera accesible sin trámites y sin coste. Al ir todo en un contenedor
Docker, el proveedor es intercambiable: la misma imagen se puede llevar a otro sitio sin tocar el
código.

\textbf{¿Es accesible la aplicación?}\\
No lo suficiente, y es la carencia que señalo en el trabajo futuro. Depende mucho del color, tanto
en los grupos anatómicos como en las escalas de métricas, y no he comprobado cómo se ve con
daltonismo. Lo primero sería añadir paletas pensadas para eso y revisar el manejo con teclado.

\textbf{¿Qué harías distinto si empezaras de nuevo?}\\
Preparar el despliegue real desde el principio. Dejarlo para el final hizo que arrastrara durante
tres sprints un montaje que servía para desarrollar pero no para publicar, y que la diferencia entre
las dos cosas no apareciera hasta que ya no quedaba margen.

\vspace{1.2em}
\accion{Si te preguntan algo que no sabes, dilo. «No lo comprobé» o «eso no lo he medido» es una
respuesta perfectamente aceptable, y es mucho mejor que inventar. Puedes añadir cómo lo
comprobarías.}

\end{document}
